\section{Exécutif national}
\label{executif-national}

\subsection{Aperçu de l'exécutif national}
\label{executif-national-apercu-de-l-executif-national}

\begin{enumerate}
 \item L'exécutif national est la direction élue de la FCEG et se compose des cinq postes suivants :
  \begin{enumerate}
   \item Présidence
   \item Vice-présidence académique
   \item Vice-présidence aux finances
   \item Vice-présidence externe
   \item Vice-présidence aux services
   \item Vice-présidence aux communications
  \end{enumerate}
 \item Les tâches et responsabilités générales de l'exécutif national incluent :
  \begin{enumerate}
   \item Faire respecter la constitution, les statuts et les politiques de la FCEG
   \item Organiser et diriger les efforts de la FCEG
   \item Assurer la réalisation de tous les mandats
   \item Gérer les personnes commissaires au sein de leurs portefeuilles respectifs, et assumer les responsabilités dans l'éventualité où elles ne sont pas complétées
   \item Établir le contact avec les instances gouvernementales lorsque nécessaire
   \item Embaucher les groupes de travail pour les personnes qui leur succéderont selon les besoins d'ici le 30 avril
   \item Créer un document de transition avant le 30 avril à conserver dans les archives électroniques de la FCEG
  \end{enumerate}
\end{enumerate}

\subsection{Présidence}
\label{executif-national-presidence}

\begin{enumerate}
 \item Les tâches et responsabilités de la présidence incluent :
  \begin{enumerate}
   \item Assurer le leadership et la supervision de l'exécutif national
   \item S'assurer que les mandats de chaque vice-présidence sont complétés
   \item Organiser les réunions de l'exécutif national
   \item S'assurer que la planification à long terme de la FCEG est entreprise, avec des mises à jour annuelles données aux membres
   \item Maintenir les canaux de communication avec toutes les parties prenantes de la FCEG
   \item Assurer la communication avec Ingénieurs Canada et assister à l'assemblée générale annuelle et aux réunions du conseil d'administration d'Ingénieurs Canada, si possible
  \end{enumerate}
\end{enumerate}

\subsection{Vice-présidence académique (VPA)}
\label{executif-national-vice-presidence-academique-vpa}

\begin{enumerate}
 \item Les tâches et responsabilités de la VPA incluent :
  \begin{enumerate}
   \item Développer une compréhension des initiatives académiques, de l'expérience de travail intégrée, du contenu des programmes d'études, des normes et exigences d'agrément ainsi que de tout changement qui pourrait affecter la formation en génie
   \item Déterminer les besoins académiques, moraux, intellectuels, culturels, sociaux et économiques de la population étudiante en génie au Canada, travailler vers des changements si nécessaire, et offrir un effet de levier national dans les questions de négociation entre la représentation étudiante et l'administration au sein des établissements des membres et d'autres organisations pertinentes
   \item Communiquer les informations pertinentes et les enjeux liés à la formation en génie ou qui affectent la population étudiante en génie au Canada aux membres et à Doyens et doyennes d'ingénierie Canada
   \item Faciliter l'échange d'informations académiques pertinentes et les efforts de revendication entre les membres et les organisations éducatives nationales
   \item Assister aux réunions du Bureau canadien d'agrément des programmes de génie et de Doyens et doyennes d'ingénierie Canada et fournir les points saillants des réunions au CA
   \item Gérer la revendication en lien avec le document de positions
   \item Diriger le groupe de travail sur la revendication, qui soutiendra le portefeuille de la VPA par des tâches de revendication incluant, mais sans s'y limiter, la recherche, la création de ressources et la rédaction de déclarations sur des sujets liés aux positions de la FCEG ou à d'autres préoccupations des membres.
   \item Fournir un soutien aux efforts de revendication des membres lorsque ceux-ci s'alignent sur les efforts de la FCEG, incluant, mais sans s'y limiter, des lettres d'appui.
  \end{enumerate}
\end{enumerate}

\subsection{Vice-présidence aux finances (VPF)}
\label{executif-national-vice-presidence-aux-finances-vpf}

\begin{enumerate}
 \item Les tâches et responsabilités de la VPF incluent :
  \begin{enumerate}
   \item Agir à titre de personne trésorière de la FCEG en :
    \begin{enumerate}
     \item Émettant les paiements sortants
     \item Déposant les fonds entrants
     \item Assurant le suivi de tous les actifs financiers
     \item Conservant un registre complet et exact des dossiers financiers de toutes les années précédentes
     \item Créant et en maintenant les états financiers actuels et historiques (état des résultats, bilan et état des flux de trésorerie)
     \item S'assurant qu'un audit annuel des finances de la FCEG est complété
    \end{enumerate}
   \item Assurer le changement en temps opportun des personnes signataires autorisées sur les comptes de la FCEG à la fin de l'exercice financier.
   \item Préparer une liste des informations de perception des frais d'adhésion pour l'année en cours et l'année précédente
   \item Recueillir le nombre de personnes étudiantes représentées par chaque membre chaque année
   \item Soumettre un budget détaillé proposé pour l'année à venir au CA
   \item Documenter les pièces justificatives pour les dépenses et les revenus, y compris les reçus, les factures, les procès-verbaux approuvant les dépenses, et tout autre document nécessaire.
  \end{enumerate}
\end{enumerate}

% Numéroté 2.4.2 comme dans le document Google, où la clause ci-dessus est 2.4.1.
\setcounter{subsubsection}{1}
\subsubsection{Vice-présidence aux finances sortante}
\label{vice-presidence-aux-finances-vpf-vice-presidence-aux-finances-sortante}

\begin{enumerate}
 \item Les tâches et responsabilités de la VPF sortante incluent :
  \begin{enumerate}
   \item Travailler aux côtés de la VPF pendant les premiers mois de son mandat
   \item Finaliser les états financiers pour l'exercice financier de son mandat
   \item Assurer l'achèvement de l'audit pour l'exercice financier de son mandat.
  \end{enumerate}
\end{enumerate}

\subsection{Vice-présidence externe (VPE)}
\label{executif-national-vice-presidence-externe-vpe}

\begin{enumerate}
 \item Les tâches et responsabilités de la VPE incluent :
  \begin{enumerate}
   \item Agir à titre de contact principal de la FCEG auprès des partenaires nationaux et internationaux existants et potentiels
   \item Explorer de nouveaux projets et des idées de coopération avec les partenaires nationaux et internationaux existants et potentiels
   \item S'assurer que tous les partenaires et partenaires en développement sont invités à tous les événements pertinents
   \item Promouvoir les opportunités internationales et les avantages de la coopération internationale auprès des membres
   \item Fournir des mises à jour régulières sur les projets pertinents, les initiatives et le travail en cours des partenaires et les résultats de ces partenariats
  \end{enumerate}
\end{enumerate}

\subsection{Vice-présidence aux services (VPS)}
\label{executif-national-vice-presidence-aux-services-vps}

\begin{enumerate}
 \item Les tâches et responsabilités de la VPS incluent :
  \begin{enumerate}
   \item Être principalement responsable de toutes les activités de la FCEG et superviser les personnes responsables d'activités.
   \item Planifier, programmer et animer des réunions de suivi mensuelles avec les personnes responsables d'activités.
   \item S'assurer que les responsabilités décrites dans les ententes d'activités sont complétées.
   \item Organiser les rapports de transition, les réunions de transition avec l'équipe sortante et entrante, et la clôture financière des activités.
   \item Encourager les membres à soumettre des candidatures et à accueillir des activités
   \item Soutenir les membres dans leurs candidatures pour accueillir des activités
   \item Recueillir les commentaires des personnes déléguées à la suite des activités dans le mois suivant la conclusion de l'activité.
  \end{enumerate}
\end{enumerate}

\subsection{Vice-présidence aux communications (VPC)}
\label{executif-national-vice-presidence-aux-communications-vpc}

\begin{enumerate}
 \item Les tâches et responsabilités de la VPC incluent :
  \begin{enumerate}
   \item Maintenir et mettre à jour des archives de fichiers originaux relatifs au logo, aux bannières, aux articles de marque et aux modèles de la FCEG.
   \item Maintenir et faire respecter les lignes directrices sur le logo et l'image de marque de la FCEG
   \item Aider les activités dans la création de modèles de conception pour les documents ou le matériel promotionnel, selon les besoins
   \item Obtenir des devis pour des articles de marque pour la FCEG ou ses activités selon les besoins
   \item Mettre à jour les profils de médias sociaux pour la FCEG
   \item Diriger le groupe de travail sur les médias et le marketing, qui créera des graphiques, recherchera les tendances et proposera de nouvelles idées sur la meilleure façon d'utiliser les réseaux sociaux de la FCEG.
   \item Diriger le groupe de travail sur le bilinguisme, qui soutiendra la traduction, élaborera des plans pour améliorer le bilinguisme lors des activités ou des réunions en ligne, et aidera à créer des ressources sur le bilinguisme pour les autres.
   \item Faciliter l'engagement entre les membres tout au long de l'année.
   \item Envoyer des publications aux membres tous les deux mois pour les tenir au courant des progrès de l'équipe de la FCEG.
  \end{enumerate}
\end{enumerate}
